\chapter{文獻探討}\label{chap:literature}
\begingroup
\setlength{\emergencystretch}{2em}
\clubpenalty=10000
\widowpenalty=10000

本章依研究問題整理相關方法。先說明具有時間資訊的分類與類別不平衡，再介紹本研究使用的分類器、集成方式及特徵選取，最後討論評估與統計比較。文獻用來交代方法的出發點與適用條件；本研究實際採用的資料分區、參數與操作順序另於第三章列明。

\section{時間資訊與分類資料的變化}

\subsection{歷史資料與近期資料的取捨}

Gama 等人\autocite{gama2014drift}將概念漂移的討論放在資料持續到達、輸入與目標關係可能隨時間改變的情境。這項觀點提醒研究者，增加歷史資料雖可能增加樣本量，卻不能直接保證模型更適合目前的資料。對本研究而言，Old 與 New 的區分提供一種觀察方式：分別建立模型，再比較保留歷史規則、使用近期規則或整合兩者的結果。

不過，年份不同不等於已證實概念漂移。某年度分數下降，可能與正類比例、樣本組成、門檻或可用紀錄數改變有關。若沒有額外的分布或漂移檢測，就無法由分數變化判定是哪一種原因。本研究因此使用「不同期間資料」描述實驗條件，將年度結果視為補充分析，不把所有差異歸因於景氣、制度或特定漂移機制。

\subsection{固定時間切分與滑動視窗}

具有時間資訊的實驗，至少需要說清楚訓練、驗證與測試資料各自來自何時。固定時間切分在選定期間建立模型，再評估保留的較晚資料；滑動視窗則會隨評估時間向前移動可使用的資料區間。兩者對應不同問題：前者觀察固定模型在後續期間的表現，後者還包含更新時間與視窗長度的影響。若使用同一組測試年份，兩者仍不能視為相同實驗。

本研究主要事件標籤的 Baseline、集成與 FS 都固定一組年份配置，藉由 Old 模型池加入 New 模型，模擬一次模型層級的更新，再以後續期間評估；Test 期間未再更新。第三章的滑動概念圖用於說明多輪配置與本次有限兩階段配置的差別，不將未執行的長序列更新寫成成果。

在資料定義上，也要區分觀測年度與事件時間。Lombardo 等人\autocite{lombardo2022bankruptcy}提供美國企業破產研究的公開資料與基準。本研究沿用其公開資料作為案例，但是否能回答提前預測問題，仍須依實際版本的標籤與可用日期判斷。本研究的公開資料只記錄公司最終狀態，第三章因此依最後觀測年度重建事件標籤，並保留最終狀態標籤作為對照。

\subsection{集成式持續學習與兩階段更新}

持續學習的更新可以發生在模型集成層級。Elwell 與 Polikar\autocite{elwell2011incremental} 的 Learn++.NSE 隨新資料批次訓練分類器，再以加權整合運用既有與新模型，提供批次增量集成的例子。此處引用其更新形式作為方法定位依據；本研究未重現其特定權重機制，也不沿用其效能結論。

本研究將此問題限縮為兩個時間階段：先以歷史資料建立 Old 模型，再以近期資料建立 New 模型，在保留 Old 候選的情況下擴充模型池，並決定如何整合。更新的是集成成員及其使用方式，不要求每個基礎分類器都延續舊模型權重。固定後續測試期用來評估這次更新的結果；它限制了可觀察的更新次數，並不取消 Old 到 Old／New 的兩階段設計。對多輪更新的遺忘、知識累積及長期穩定性，仍需其他評估。

\section{類別不平衡與採樣方法}

\subsection{不平衡問題的觀察方式}

不平衡資料的困難不只取決於正負類數量。Batista 等人\autocite{batista2004balancing}比較多種平衡方法，指出少數類樣本不足與類別重疊等因素也會影響學習。因此，調整類別比例並不是唯一目的；採樣可能同時改變類別邊界附近的資料與模型接觸到的樣本資訊。

本研究依此將不採樣、低採樣、過採樣與混合採樣列為不同比較條件。不採樣保留原始分布，讓其他設定具有共同基準。採樣只作用於允許的訓練資料，驗證與測試仍保留原有比例，才能觀察模型在未經人為平衡的資料上有何表現。

\subsection{低採樣與 Tomek links}

低採樣藉由移除部分樣本改變訓練資料，但不同方法的移除規則並不相同。Tomek links 關注彼此接近、卻屬於不同類別的樣本配對，依設定移除其中指定類別的樣本。相較於隨機刪除大量多數類，這類方法具有清理類別邊界的用意。官方實作也明確指出，移除 Tomek links 後不保證各類樣本數相同\autocite{imblearn-tomek-docs}。

因此，本研究表中的 Under 代表具體演算法 TomekLinks，不能解讀為每次都將正負類調成一比一。此點也影響 Old／New 模型池的理解：三種採樣來源提供的是不同資料處理方式，不只是三種固定類別比例。

\subsection{過採樣：SMOTE 與 ADASYN}

Chawla 等人\autocite{chawla2002smote}提出 SMOTE，利用少數類樣本與其鄰近樣本之間的資訊建立合成樣本，擴充少數類的訓練表示。它與單純複製既有樣本不同，但合成資料仍由原訓練樣本推得，不能當成額外蒐集的獨立觀測。

He 等人\autocite{he2008adasyn}提出 ADASYN，依少數類樣本周邊的學習困難程度分配合成數量，使較難處理的區域獲得較多合成樣本。本研究的 Over 使用 ADASYN。這項設計使過採樣不只增加正類數量，也會依原資料的鄰近關係改變合成樣本的分布。

由於這些方法依賴距離或鄰近樣本，特徵尺度與選取結果都可能影響採樣。若先在全部資料上採樣再切分，合成程序便可能間接使用評估資料的資訊。前處理與採樣皆限制於允許的擬合分區；FS 實驗則先選定特徵，再合成樣本。

\subsection{混合採樣與 SMOTEENN}

混合採樣將增加樣本與清理樣本結合。Batista 等人\autocite{batista2004balancing}討論 SMOTE 搭配 Tomek links 或 Edited Nearest Neighbours（ENN）的做法，嘗試在擴充少數類之後整理重疊區域。本研究使用 SMOTEENN 作為 Hybrid：先由 SMOTE 合成樣本，再依 ENN 規則清理。

表~\ref{tab:lit-sampling}整理本研究採樣設定的作用差異。三種方法的目的不同，本文不預設某一方法普遍最佳，而是讓它們分別形成基礎模型，進一步比較模型整合後的結果。

\begin{table}[htbp]
\centering\small
\setstretch{1.15}\setlength{\tabcolsep}{4pt}
\caption{本研究採樣設定與方法作用}\label{tab:lit-sampling}
\begin{tabular}{@{}>{\raggedright\arraybackslash}p{0.16\linewidth}>{\raggedright\arraybackslash}p{0.26\linewidth}>{\raggedright\arraybackslash}p{0.48\linewidth}@{}}
\toprule
設定 & 代表方法 & 主要作用與解讀注意事項 \\
\midrule
None & 不採樣 & 保留訓練分布，作為採樣比較的基準。 \\[4pt]
Under & TomekLinks & 清理跨類別鄰近配對；不保證類別數相等。 \\[4pt]
Over & ADASYN & 依少數類周邊的困難程度分配合成樣本。 \\[4pt]
Hybrid & SMOTEENN & 結合少數類合成與鄰近規則清理。 \\
\bottomrule
\end{tabular}
\par\smallskip{\footnotesize (資料來源：本研究依 SMOTE、ADASYN、混合採樣文獻與 TomekLinks 文件整理\autocite{chawla2002smote,he2008adasyn,batista2004balancing,imblearn-tomek-docs})\par}
\end{table}

\section{表格資料分類器}

\subsection{線性模型與樹模型}

本研究的 Baseline 包含不同模型家族。Logistic Regression 以線性分數連結二元類別機率，適合提供結構較簡單的比較基準；Random Forest 及梯度提升樹則可以透過樹的分割表示較複雜的關係。比較這些模型，有助於觀察資料期間與採樣的差異是否只出現在某一種架構。

Breiman\autocite{breiman2001rf}提出的 Random Forest 結合多棵樹，並在建立樹時引入隨機性；其表現與個別樹的能力及樹之間的相關程度有關。XGBoost 則是梯度提升樹的系統實作，Chen 與 Guestrin\autocite{chen2016xgboost}將正則化目標與可擴充的樹學習設計結合。Ke 等人\autocite{ke2017lightgbm}提出 LightGBM，著重提升梯度提升樹處理資料的效率。這些方法雖都包含多棵樹，其訓練方式與超參數作用並不相同。

本研究在用語上區分兩個層次：Random Forest、XGBoost 或 LightGBM 本身可由多棵樹構成；第三章的 Old／New 集成則整合多個分別訓練的分類器。後者的成員差異來自資料期間與採樣方式，不能直接與模型內部的樹數混為一談。研究框架不預設某一分類器，先由共同驗證準則選出最佳基礎模型及參數，再比較其集成規則；目前測試暫用 XGBoost 檢查流程，不能將此測試設定當成最佳模型的選擇結論。

\subsection{表格深度學習與情境學習}

在表格深度學習方面，Gorishniy 等人\autocite{gorishniy2025tabm}提出 TabM，藉由共享大部分參數，在單一模型內產生多個預測，結合多層感知器與有效率的集成設計。TabM 的模型內部集成與本研究 Old／New 模型池的整合屬於不同層次。本文將 TabM 保留為深度學習 Baseline，第三章補列已保存設定，第四章先呈現既有最終狀態標籤結果；事件標籤分數待重跑後更新。

Qu 等人\autocite{qu2025tabicl}提出 TabICL，採用情境學習（in-context learning），以提供的資料作為預測情境，而非必須針對每份目標資料更新全部模型參數。該文介紹的是原方法的架構與基準結果，不能直接替本研究歷史執行紀錄背書。

本研究已完成事件標籤實驗的四種傳統分類器，在每個訓練策略與採樣組合下均評估 24 組候選參數；候選數相同，但搜尋空間與計算成本不同，不能解釋為各模型均已充分調參。TabM 與 TabICL 的既有紀錄另列於第四章的暫列結果表，並保留為六分類器架構的兩個深度學習分支。由於其標籤與驗證方式不同，暫列數值不參與事件標籤的同條件排名；後端及執行範圍的限制於附錄說明。

\section{靜態集成與動態選擇}

\subsection{靜態成員與機率整合}

靜態集成在推論前固定參與的分類器與整合規則。若各成員輸出正類機率，可以用等權平均形成集成分數。此作法容易說明，但並不保證加入越多模型越好：若某些成員在相同樣本上給出不合適的分數，平均也可能降低區辨能力。

本研究從相同六模型池比較數種靜態安排。Old3 與 New3 各整合同一期間的模型，Equal6 使用全部成員；V-static 則由驗證資料選擇同時含 Old 與 New 的混合子集。這些控制組對應不同問題：是否需要跨期間、是否應納入全部成員，以及驗證選擇子集是否有幫助。選模資料與測試資料必須分開，否則將測試分數最高的子集當成可使用的方法，會高估實際選擇程序的能力。

\subsection{動態分類器選擇}

動態分類器選擇（Dynamic Classifier Selection, DCS）依待預測樣本的局部區域選擇分類器。Woods 等人\autocite{woods1997local}以局部分類正確程度研究模型組合，使選擇與樣本周邊的表現連結。本研究比較 Overall Local Accuracy（OLA）與 Local Class Accuracy（LCA）。OLA 使用鄰近區域的整體正確程度；LCA 則依分類器對查詢樣本所預測的類別，查看相應鄰近樣本的表現\autocite{deslib-ola-docs,deslib-lca-docs}。

局部能力並不是在測試標籤上計算，而是由另行保留的 Dynamic Selection Dataset（DSEL）估計。DSEL 因而是動態選擇規則的參考資料，不是可再併回基礎模型訓練的閒置資料。若這個區域不足以代表待預測樣本，選出的模型也未必適合，因此 DCS 的效果仍須透過保留測試資料確認。

\subsection{動態集成選擇}

動態集成選擇（Dynamic Ensemble Selection, DES）將選擇對象由單一分類器擴大為一組分類器。Ko 等人\autocite{ko2008dynamic}討論由動態分類器選擇延伸至動態集成選擇的方式。本研究比較 KNORA-E 與 KNORA-U：前者找出能正確處理局部鄰近樣本的成員，必要時縮小區域；後者利用能正確處理部分鄰近樣本的成員及其局部表現進行整合\autocite{deslib-knorae-docs,deslib-knorau-docs}。

DESlib 提供靜態、DCS 與 DES 的實作與共同介面\autocite{cruz2020deslib}。本研究依固定版本設定使用該套件，讓動態方法共用相同模型池、DSEL 與近鄰數，並在第三章交代機率整合方式。此處只保留理解選擇機制所需的原理；具體選擇、權重與例外處理由實際使用的演算法決定。

局部區域依特徵距離建立，因此不同尺度或特徵子集可能改變鄰居，即使分類器種類不變，動態選擇結果也可能不同。這是本研究將 FS 與集成共同評估的理由之一。若 FS 後分數改變，應考慮基礎模型與局部區域同時改變，而非只用「移除雜訊」解釋所有結果。

\section{特徵選取與可解釋方法}

\subsection{Filter、Wrapper 與 Embedded}

Guyon 與 Elisseeff\autocite{guyon2003selection}整理特徵選取的多種問題，包括預測能力、計算成本及資料理解。依特徵評估與模型訓練的關係，常區分為過濾式（Filter）、包裝式（Wrapper）及嵌入式（Embedded）方法。這項分類說明方法如何產生特徵子集，並不直接表示哪一類在所有資料上較佳。

本研究以互資訊（Mutual Information, MI）作為 Filter 代表，依個別特徵與標籤的關聯排序；以遞迴特徵刪除（Recursive Feature Elimination, RFE）作為 Wrapper 代表，反覆依估計器的特徵重要性移除特徵；另以決策樹（CART）訓練後的重要性作為 Embedded 代表。各方法在特徵間關係、模型依賴與搜尋成本上有所不同，實際操作可由所用套件的特徵選取介面確認\autocite{sklearn-fs-docs}。

三類方法在本研究是平行條件，並未串接成先 Filter、再 Wrapper、最後 Embedded 的程序。各方法分別產生固定特徵集合，再重新建立模型池。保留 50\% 與 80\% 是事先設定的兩種精簡程度，用來比較刪減較多與保留較多資訊的差別，不稱為文獻認定或本資料搜尋得到的最佳比例。

\subsection{SHAP 的角色與特徵排序}

可解釋人工智慧（Explainable Artificial Intelligence, XAI）關心如何說明模型的判斷；特徵選取則決定模型之後使用哪些輸入，兩者目的不同。Lundberg 與 Lee\autocite{lundberg2017shap}提出 SHAP，以特徵貢獻值說明個別預測。將多筆樣本的貢獻取絕對值後平均，可以形成整體重要性排序，再依保留數篩選特徵。

Hancock 等人\autocite{hancock2025shapfs}研究利用 SHAP 進行特徵選取與分析，提供將解釋資訊用於篩選的例子。本研究將 SHAP 視為另一種特徵排序方式，與 MI、CART 及 RFE 比較，並未將它稱為最新 FS，也未直接重現該文的完整選取程序。實作使用 LightGBM 的樹模型貢獻輸出，排除偏置項後計算平均絕對貢獻；模型、樣本及保留數均在第三章說明。

SHAP 值說明的是特定模型如何使用特徵，並不證明特徵對公司狀態具有因果影響。若兩個財務欄位高度相關，重要性也可能分配在不同欄位，不能只依一次排序認定其中之一沒有資訊。Goldwasser 與 Hooker\autocite{goldwasser2025rankings}進一步研究特徵重要性排名的統計顯著性，提醒排名亦有不確定性。本研究附錄只整理十個種子中的入選次數，屬於描述性紀錄，沒有執行該文的排名檢定。

\subsection{特徵選取的驗證界線}

只要選取程序使用了標籤或資料分布，它就屬於模型建立的一部分。如果先用測試資料選出特徵，再評估模型，分數便包含了測試資訊。本研究讓各 FS 只在 Old 擬合資料上執行，將同一組欄位及順序套用至 New、DSEL、Calibration 與 Test。如此可比較固定歷史選取規則在後續資料的表現，也同時限制了模型無法使用近期資料重新挑選特徵。

另外，FS 在採樣之前執行，讓選取依據維持原擬合資料的比例。這是一項固定流程的設計選擇，而非宣稱所有資料都應依此順序處理。不同順序可能得到不同子集，本研究並未另做先採樣再選取的對照，因此結論限於目前程序。

\section{分類評估與統計比較}

\subsection{排序能力與指定門檻的結果}

AUC 衡量模型將正類分數排在負類之前的能力，不依賴單一分類門檻。F1-score 與 Recall 則需要先把分數轉成類別。兩者回答的問題不同：分數排序改善時，固定門檻下仍可能漏掉部分正類；提高 Recall 也可能伴隨更多誤報。因此，本研究以 AUC 為主要報告指標，並同時列出 F1-score 與 Recall。

Precision--Recall 關係較直接呈現正類預測的精確程度\autocite{saito2015pr}。本研究另以 Average Precision（AP）比較驗證資料中的靜態候選。AP 依召回增量加權精確率，與以梯形積分計算的 PR 曲線面積不應直接視為相同數值\autocite{sklearn-ap-docs}。本文選成員使用 Calibration AP，選分類門檻使用 Calibration F1-score，最終 Test 的三項指標各自保留，不用測試結果回頭調整上述規則。

\subsection{重複紀錄、不確定性與多重比較}

Efron 與 Tibshirani\autocite{efron1986bootstrap}說明以重抽樣估計統計量變異的基本原理；但重抽樣單位必須符合資料結構。Obuchowski\autocite{obuchowski1997clustered}討論群聚 ROC 資料的相關性問題。對本研究而言，同一公司可能有數年紀錄，若逐列當成獨立樣本重抽，便忽略公司內部的關聯。因此，第四章補充分析以公司為單位，讓同公司年度紀錄一起被抽入，兩方法也使用相同抽樣公司，形成配對差值。

隨機種子的標準差與公司重抽樣信賴區間也有不同用途。前者表示既定資料下執行隨機程序的分數變化；後者表示已保存模型在重抽測試公司下的差值變化。兩者都不能完整代表重新取得訓練資料、變更年份或換到另一市場時的不確定性。

當同時比較多種方法與指標，單獨對每項使用 0.05 門檻會累積誤判機會。Holm\autocite{holm1979multiple}提出逐步多重檢定校正。本研究對事件標籤與最終狀態標籤分別進行事後分析，各將兩項實驗的 165 項比較放在同一校正範圍；兩種標籤的校正家族分開，完整保留增加、下降及未顯著的結果；逐項信賴區間與校正後的顯著判斷分別呈現，不把未顯著解釋為兩方法已證實相同。

\section{相關研究整理與本研究定位}

表~\ref{tab:lit-position}依方法角色整理代表性研究。本表比較的是研究內容與本研究使用方式，不是把不同資料集、切分與指標下的分數放在一起排名。

\begingroup\small\setstretch{1.12}\setlength{\tabcolsep}{4pt}
\begin{longtable}{@{}>{\raggedright\arraybackslash}p{0.24\linewidth}>{\raggedright\arraybackslash}p{0.30\linewidth}>{\raggedright\arraybackslash}p{0.36\linewidth}@{}}
\caption{相關研究主題與本研究的銜接}\label{tab:lit-position}\\
\toprule
研究 & 討論重點 & 本研究的使用方式與界線 \\
\midrule\endfirsthead
\multicolumn{3}{c}{表~\thetable（續）}\\\toprule
研究 & 討論重點 & 本研究的使用方式與界線 \\
\midrule\endhead
\midrule\multicolumn{3}{r}{續下頁}\\\endfoot
\bottomrule\endlastfoot
Chawla 等人（2002）\autocite{chawla2002smote} & 少數類合成樣本 & 作為混合採樣中 SMOTE 的原理依據；合成資料不是獨立新觀測。 \\[4pt]
He 等人（2008）\autocite{he2008adasyn} & 依困難程度分配合成數量 & Over 採 ADASYN；實際影響依本資料比較，不沿用原文成效。 \\[4pt]
Gama 等人（2014）\autocite{gama2014drift} & 時間變化與適應問題 & 說明歷史／近期資料取捨；本研究未執行線上漂移偵測。 \\[4pt]
Woods 等人（1997）與 Ko 等人（2008）\autocite{woods1997local,ko2008dynamic} & 局部能力與動態模型選擇 & 比較 DCS／DES，使用獨立 DSEL；不假設一定優於靜態方法。 \\[4pt]
Lombardo 等人（2022）\autocite{lombardo2022bankruptcy} & 企業資料與分類基準 & 提供資料研究背景；依實際公開版另核對標籤與識別碼數。 \\[4pt]
Hancock 等人（2025）\autocite{hancock2025shapfs} & SHAP 用於特徵選取與分析 & 將 SHAP 排序列為 FS 對照，未重現其全部程序。 \\[4pt]
Goldwasser 與 Hooker（2025）\autocite{goldwasser2025rankings} & 特徵重要性排名的不確定性 & 入選次數僅作描述，不冒稱已完成排名顯著性檢定。 \\*
\end{longtable}
\noindent{\footnotesize (資料來源：本研究依表列原始文獻整理)}
\endgroup

上述文獻分別說明採樣、時間變化、動態選擇與 FS 的原理，但這些原理本身不能回答本資料中的歷史模型是否仍有互補資訊，也不能保證縮減特徵後的局部選擇更可靠。本研究將待回答的實證問題具體化為三項對照：以 Old、New 與 Retrain 觀察資料期間的差異；讓靜態與動態方法共用六模型池、DSEL 及驗證決策規則；再於相同測試公司上比較 FS 與 no-FS。事件標籤與最終狀態標籤另行對照，以檢查方法排序是否依賴目標定義。

本研究的補充價值在於將這些問題置於可追溯的共同流程，並以公司層級的配對不確定性限制結論。此定位不代表既有文獻從未結合上述方法，也不主張新的演算法或跨資料集的普遍優越性；現有證據來自單一公開資料、固定年份與 XGBoost 的測試實作；研究框架的基礎模型仍由驗證表現選定，正式選模後須重新評估。
\endgroup
